\documentclass[../main.tex]{subfiles}
\begin{document}
\chapter{Fourier Transform}
\section{Definitions and Simple Properties}
\begin{definition}[Fourier Transform]
  For an integrable function $f: \R \to \C$, we define its \textit{Fourier Transform} as the function $\hat{f}: \R \to \C$ given by:
  \[
    \hat{f}(\lambda) = \int_{-\infty}^{\infty} e^{-i\lambda x}f(x) \d{x}
  \]
  for $\lambda \in \R$.
\end{definition}
\begin{remark}[Notation]
  We denote the linear map taking a function $f$ to its Fourier transform as:
  \[
    \mathcal{F}(f) = \hat{f}
  \]
  We also write:
  \[
    \mathcal{F}_{x \to \lambda}[f(x)] = \hat{f}(\lambda)
  \]
  to mean the Fourier transform of $f$ in the variable $x$ evaluated at $\lambda$.
\end{remark}
We shall assume that $f$ and its derivatives decay sufficiently fast so that we can integrate by parts.

Using the Riemann-Lebesgue lemma (\cref{RLL}), we have:
\[
  \hat{f}(\lambda) \to 0 \text{ as } |\lambda| \to \infty
\]
\begin{proposition}
  For integrable $f: \R \to \C$, we have the following: \label{fourierProps}
  \begin{enumerate}
    \item Fourier transforms of derivatives introduce powers of $\lambda$:
      \[
        \mathcal{F}_{x \to \lambda}\left[\left(\deriv{}{x}\right)^{k} f(x)\right] = (i \lambda)^{k} \hat{f}(\lambda)
      \]
    \item Dual to \textbf{(i)}, introducing powers of $x$, introduces derivatives into the Fourier transform:
      \[
        \mathcal{F}_{x \to \lambda}\left[x^k f(x)\right] =
      \]
    \item Shifting the function introduces an exponential into the Fourier transform:
      \[
        \mathcal{F}_{x \to \lambda}\left[f(x - a)\right] = e^{-i \lambda a} \hat{f}(\lambda)
      \]
    \item Dual to \textbf{(iii)}, introducing an exponential shifts the Fourier series:
      \[
        \mathcal{F}_{x \to \lambda}\left[e^{-iax}f(x)\right] = \hat{f}(\lambda + a)
      \]
  \end{enumerate}
\end{proposition}
\begin{proof}[\textbf{i}]
  If we integrate by parts $k$ times, then because the exponential term disappears at $\pm \infty$, the boundary terms disappear and so we obtain:
  \begin{align*}
    \mathcal{F}_{x \to \lambda}\left[\left(\deriv{}{x}\right)^{k} f(x)\right] &= \int_{-\infty}^{\infty} f(x) \left(-\deriv{}{x}\right)^{k} e^{-i\lambda x} \d{x} \\
                                                                              &= \int_{-\infty}^{\infty}  f(x) (i \lambda)^{k} e^{-i\lambda x} \d{x} \\
                                                                              &= (i \lambda)^{k} \hat{f}(\lambda)
  \end{align*}
  as required.
\end{proof}
\begin{proof}[\textbf{ii}]
  In this case, it easier to start from the right hand side of the expression.
  Differentiating under the integral sign, we have:
  \begin{align*}
    \left(i \deriv{}{\lambda}\right)^{k} \hat{f}(\lambda) &= \left(i \deriv{}{\lambda}\right)^{k} \int_{-\infty}^{\infty} e^{-i \lambda x}f(x) \d{x} \\
                                                          &= \int_{-\infty}^{\infty} f(x) i^k \deriv[k]{}{\lambda}(e^{-i\lambda x}) \d{x}\\
                                                          &= \int_{-\infty}^{\infty} -i^{2k} x^k e^{-i \lambda x} f(x) \d{x} \\
                                                          &= \int_{-\infty}^{\infty} e^{-i\lambda x} x^{k} f(x) \d{x} \\
                                                          &= \mathcal{F}_{x \to \lambda}[x^{k} f(x)]
  \end{align*}
  as required.
\end{proof}
\begin{proof}[\textbf{iii}]
  We see that:
  \begin{align*}
    \mathcal{F}_{x \to \lambda}\left[f(x - a)\right] &= \int_{-\infty}^{\infty} e^{-i\lambda x} f(x - a) \d{x} \\
                                                     &= \int_{-\infty}^{\infty} e^{-i \lambda (x + a)} f(x) \d{x} \text{ substituting $x \mapsto x + a$} \\
                                                     &= e^{-i \lambda a} \int_{-\infty}^{\infty} e^{-i\lambda x} f(x) \d{x} \\
                                                     &= e^{-i \lambda a} \hat{f}(\lambda)
  \end{align*}
  as required.
\end{proof}
\begin{proof}[\textbf{iv}]
  We see that:
  \begin{align*}
    \mathcal{F}_{x \to \lambda}\left[e^{-iax}f(x)\right] &= \int_{-\infty}^{\infty} e^{-i\lambda x} e^{-iax} f(x) \d{x} \\
                                                         &= \int_{-\infty}^{\infty} e^{-i(\lambda + a)x} f(x) \d{x} \\
                                                         &= \hat{f}(\lambda + a)
  \end{align*}
\end{proof}
Property \textbf{(i)} above turns out to be very useful for solving ODEs.
Suppose we are trying to solve the problem:
\begin{equation}
  p\left(\deriv{}{x}\right)y = f \text{ for } x \in \R \label{polyProb}
\end{equation}
where $p$ is a polynomial expression evaluated on the derivative operator, giving another derivative operator, i.e.:
\[
  p(z) = \sum_{i = 0}^{n} a_n z^n \implies p\left(\deriv{}{x}\right) = \sum_{i = 0}^{n} a_i \deriv[i]{}{x}
\]
Taking the Fourier transform of one of these terms, using \cref{fourierProps} \textbf{(i)}, we have:
\[
  \mathcal{F}_{x \to \lambda}\left[a_i \deriv[i]{}{x} y\right] = a_i(i\lambda)^i \hat{y}(\lambda)
\]
and so taking Fourier transforms of both sides of \cref{polyProb}:
\[
  P(i \lambda) \hat{y}(\lambda) = \hat{f}(\lambda) \text{ for } \lambda \in \R
\]
Thus, provided $P(i \lambda) \neq 0$, we can obtain the Fourier transform of the solution $y$:
\[
  \hat{y}(\lambda) = \frac{\hat{f}(\lambda)}{P(i\lambda)}
\]
It would be useful if we could invert the Fourier transform as this would allow us to recover $y$ from the expression for its Fourier transform $\hat{y}$ above.
It turns out that we can in fact do this.
\begin{proposition}[Inverting Fourier transforms]
  The Fourier transform $\hat{f}$ of an integrable function $f: \R \to \C$ can be inverted using the formula: \label{fourierInverse}
  \[
    f(x) = \frac{1}{2\pi} \int_{-\infty}^{\infty} e^{i \lambda x} \hat{f}(\lambda) \d{\lambda} \text{ for } x \in \R
  \]
\end{proposition}
\begin{proof}
  Consider the integral:
  \begin{align*}
    I_n(x) &= \frac{1}{2\pi} \int_{-n}^{n} e^{i \lambda x} \hat{f}(\lambda) \d{\lambda} \\
           &= \frac{1}{2\pi} \int_{-n}^{n} e^{i\lambda x} \left[\int_{-\infty}^{\infty} e^{-i\lambda y} f(y) \d{y}\right] \d{\lambda} \\
           &= \int_{-\infty}^{\infty} f(y) \left[\frac{1}{2\pi} \int_{-n}^{n} e^{i\lambda(x - y)} \d{\lambda}\right] \d{y} \\
           &= \int_{-\infty}^{\infty} f(y) \cdot \frac{\sin[n(x - y)]}{\pi(x - y)} \d{y} \\
           &= \int_{-\infty}^{\infty} \frac{\sin[ny]}{\pi y} f(x + y) \d{y} \text{ by substituting $y \mapsto x + y$}
  \end{align*}
  Recall from IA Differential Equations Example Sheet 1 Q13 that:
  \[
    \int_{-\infty}^{\infty} \frac{\sin[ny]}{\pi y} \d{y} = 1 \text{ for all } n > 0
  \]
  Thus, we can write:
  \begin{align*}
    I_n(x) - f(x) &= \int_{-\infty}^{\infty} \frac{\sin[ny]}{\pi y}[f(x + y) - f(x)] \d{y} \\
                  &= \int_{-\infty}^{\infty} \sin[ny] F(x, y) \d{y}
  \end{align*}
  where we define:
  \[
    F(x, y) = \frac{f(x + y) - f(x)}{\pi y}
  \]
  Integrating by parts, we obtain:
  \[
    I_n(x) - f(x) = \frac{1}{n} \int_{-\infty}^{\infty} \cos[ny] F_y(x, y) \d{y}
  \]
  Taking $n \to \infty$, this vanishes and so $I_n \to f(x)$, i.e.:
  \[
    f(x) = \lim_{n \to \infty} I_n(x) = \frac{1}{2\pi} \int_{-\infty}^{\infty} e^{i \lambda x} \hat{f}(\lambda) \d{\lambda}
  \]
\end{proof}
If we remove the limiting process in the above proof, then:
\[
  f(x) = \int_{-\infty}^{\infty} f(y) \left[\frac{1}{2\pi} \int_{-\infty}^{\infty} e^{i\lambda(x - y)} \d{\lambda}\right] \d{x}
\]
so it might be reasonable to conclude that:
\[
  \delta(x) = \frac{1}{2\pi} \int_{-\infty}^{\infty} e^{i\lambda x} \d{\lambda}
\]
when used inside of an appropriate integral.
We can use this to derive an result analogous to Parseval's theorem for Fourier series from \cref{parsevalSec}.
\begin{proposition}[Parseval's theorem for Fourier transform]
  For integrable functions $f, g: \R \to \C$, we have:
  \[
    \int_{-\infty}^{\infty} f(x) \overline{g(x)} \d{x} = \frac{1}{2\pi} \int_{-\infty}^{\infty} \hat{f}(\lambda) \overline{\hat{g}(\lambda)} \d{\lambda}
  \]
\end{proposition}
\begin{proof}
Starting from the right hand side, using \cref{fourierInverse}, we have:
\begin{align*}
  \frac{1}{2\pi} \int_{-\infty}^{\infty} \hat{f}(\lambda) \overline{\hat{g}(\lambda)} \d{\lambda} &= \frac{1}{2\pi} \int_{-\infty}^{\infty} \left[\int_{-\infty}^{\infty} e^{-i \lambda y} f(y) \d{y}\right]\left[\int_{-\infty}^{\infty} e^{i\lambda x} \overline{g(x)} \d{x}\right] \d{\lambda} \\
                                                                                                  &= \int_{-\infty}^{\infty} \int_{-\infty}^{\infty}  f(y) \overline{g(x)}\underbrace{\left[\frac{1}{2\pi} \int_{-\infty}^{\infty} e^{i\lambda(x - y)} \d{\lambda}\right]}_{\delta(x - y)} \d{x} \d{y} \\
                                                                                                  &= \int_{-\infty}^{\infty} f(x) \overline{g(x)} \d{x}
\end{align*}
as required.
\end{proof}
\begin{remark}
  Taking $g = f$ in the above, we obtain:
  \[
    \int_{-\infty}^{\infty} |f(x)|^2 \d{x} = \frac{1}{2\pi} \int_{-\infty}^{\infty} |\hat{f}(\lambda)|^2 \d{\lambda}
  \]
\end{remark}
\begin{definition}[Convolution]
  For integrable functions $f, g: \R \to \C$, we define their \textit{convolution} as:
  \[
    (f \ast g)(x) = \int_{-\infty}^{\infty} f(x - y)g(y) \d{y}
  \]
\end{definition}
\begin{lemma}
  For integrable functions $f, g: \R \to \C$, $f \ast g = g \ast f$.
\end{lemma}
\begin{proof}
  We have:
  \begin{align*}
    (f \ast g)(x) &= \int_{-\infty}^{\infty} f(x - y)g(y) \d{y} \\
                  &= \int_{-\infty}^{\infty} f(y)g(x - y) \d{y} \text{ by substituting $y \mapsto x - y$} \\
                  &= (g \ast f)(x)
  \end{align*}
  This is true for any $x$ and so $f \ast g = g \ast f$.
\end{proof}
\begin{proposition}[Convolution Theorem]
  For integrable functions $f, g: \R \to \C$, we have:
  \[
    \mathcal{F}[f \ast g] = \hat{f} \hat{g}
  \]
\end{proposition}
\begin{proof}
  Starting from the left hand side:
  \begin{align*}
    \mathcal{F}[f \ast g] &= \int_{-\infty}^{\infty} e^{-i \lambda x} \left[\int_{-\infty}^{\infty} f(x - y) g(y) \d{y}\right] \d{x} \\
                          &= \int_{-\infty}^{\infty} \left[\int_{-\infty}^{\infty} f(x - y) g(y) e^{-i\lambda x} \d{x}\right] \d{y} \\
                          &= \int_{-\infty}^{\infty} \left[\int_{-\infty}^{\infty} f(x) g(y) e^{-i\lambda(x + y)} \d{x}\right] \d{y} \text{ taking $x \mapsto x  + y$} \\
                          &= \left[\int_{-\infty}^{\infty} e^{-i \lambda x} f(x) \d{x}\right]\left[\int_{-\infty}^{\infty} e^{-i \lambda y} g(y) \d{y}\right] \\
                          &= \hat{f}(\lambda) \hat{g}(\lambda)
  \end{align*}
  as required.
\end{proof}
\end{document}
